\chapter{Типизированное K43-вложение минимального
\texorpdfstring{$SU(2)+8D$}{SU(2)+8D} носителя}
\label{ch:v10-discrete-resolution-su2-k43-embedding}

Условный минимальный носитель допускает точное разложение
\begin{equation}
 \mathcal K_{43}
 =\bigl(\mathbb C^2\otimes\mathbb C^{16}\bigr)
 \oplus\mathbb C^{11}.
\end{equation}
Первые $32$ направления несут шестнадцать вейлевских дублетов, то есть
восемь дираковских дублетов; остальные одиннадцать являются синглетами.
Генераторы
\begin{equation}
 T_a=\frac{\sigma_a}{2}\otimes I_{16}\oplus0_{11}
\end{equation}
удовлетворяют $[T_a,T_b]=i\epsilon_{abc}T_c$ и
\begin{equation}
 \sum_aT_a^2=\frac34P_{32},\qquad
 \Tr(T_aT_b)=8\delta_{ab}.
\end{equation}
Поэтому полный фермионный индекс равен $8$ и
\begin{equation}
 b=-\frac{22}{3}+\frac23\,8=-2.
\end{equation}
Локальный кубический anomaly tensor тождественно нулевой, а число
вейлевских дублетов $16$ чётно, поэтому Witten anomaly отсутствует.

Однако представление имеет большой коммутант размерности
$16^2+11^2=377$: само $K_{43}$ не выбирает конкретное разложение
кратностей. Более существенно, ранговый проектор внутри активного дублета
имеет ненулевой gauge defect, равный $1$. Ранговый инвариантный проектор
можно поместить в одиннадцатимерное синглетное дополнение, но тогда его
пересечение с AF-сектором равно нулю. Наименьший инвариантный проектор в
активном секторе есть целый дублет и имеет ранг $2$.

\begin{theorem}[Условное K43-вложение существует]
Носитель $SU(2)+8D$ точно помещается в $K_{43}$, имеет $b=-2$ и проходит
локальную и глобальную anomaly-проверки. Архитектура вложения закрыта.
\end{theorem}

\begin{nogocriterion}[Ранговый полюс не является активным синглетом]
В активном дублетном секторе нет $SU(2)$-инвариантного проектора ранга
один: инвариантные проекторы содержат полные дублеты и имеют чётный ранг.
Синглетный ранговый полюс лежит вне AF-сектора. Поэтому требуется
композитный gauge-singlet полюс, а не одиночная дублетная линия.
\end{nogocriterion}

\begin{protocol}\begin{enumerate}
\item условная архитектура вложения: \textbf{$12/12$};
\item точный beta-коэффициент: \textbf{$1/1$};
\item локальная и глобальная anomaly-проверки: \textbf{$2/2$};
\item активный инвариантный полюс ранга один: \textbf{$0/1$};
\item физическое происхождение: \textbf{$0/3$};
\item следующий гейт: \textbf{композитный $SU(2)$-синглетный полюс}.
\end{enumerate}\end{protocol}

Машинный сертификат сохранён в
\path{s2t_v10_cell_birth_four_volume_nonequilibrium_bath_discrete_resolution_transmuted_mode_asymptotically_free_su2_eight_dirac_k43_typed_embedding_gate_results.json}.